% This chapter is a mess--ehb

\chapter{Chapter 9: Compressed Position Report Data Formats}


In compressed data format, the Information field contains the
station’s latitude and longitude, together with course and speed or
pre-calculated radio range or altitude.

This information is compressed to minimize the length of the
transmitted packet (and therefore improve its chances of being
received correctly under less than ideal conditions).

The Information field also contains a display Symbol Code, and there
may optionally be a plain text comment (uncompressed) as well.

\section {The Advantages of Data Compression}

Compressed data format may be used in place of the numeric lat/long
coordinates already described, such as in the !, /, @ and = formats.
Data compression has several important benefits:

\begin{itemize}
  
\item Fully backwards compatible with all existing formats.

\item Fully supports any comment string.

\item Speed is accurate to +/-1 mph up to about 40 mph and within 3\%
  at 600 mph.

\item Altitude in feet is accurate to +/- 0.4\% from 1 foot to 3000 miles.

\item Consistent one-algorithm processing of compressed latitude and
  longitude.

\item Improved position to 1 foot worldwide.

\item Pre-calculated radio range, compressed to one byte.

\item Potential 50\% compression of every position format on the air.

\item Potential 40\% reduction of raw GPS NMEA data length.

\item Additional 7-byte reduction for NEMA GGA altitudes.

\item Support for TNC compression at the NMEA source (from the GPS
  receiver).

\item Digipeater compression of old NMEA trackers on the fly.

\item Usage is optional in all cases.

\end{itemize}

The only minor disadvantages are that the course only resolves to +/- 2
degrees, and this format does not support PHG.


\section {Compressed Data Format}

Compressed data may be generated in several ways:

\begin{itemize}
    
\item by APRS software.

\item pre-entered manually into a digipeater’s beacon text.

\item by a digipeater converting raw tracker NMEA packets to compressed.

\end{itemize}


[In future, there is the possibility that a Kantronics KPC-3 or other tracker
TNC will be able to compress data directly from an attached GPS receiver].
In all cases the compressed format is a fixed 13-character field:

/YYYYXXXX\$csT

where

\begin{itemize}

\item / is the Symbol Table Identifier
\item YYYY is the compressed latitude
\item XXXX is the compressed longitude
\item \$ is the Symbol Code
\item cs is
  \begin {itemize}
  \item  the compressed course/speed;
  \item compressed pre-calculated radio range;
  \item or compressed altitude
    \end {itemize}
\item T is the compression type indicator

\end {itemize}

\begin{tabular}{|c|c|c|c|c|c|}
  \hline
  \multicolumn {6}{|l|}{Compressed Position Data} \\
  \hline
  Sym & Compressed & Compressed & Symbol & Compressed   & Comp \\ [-1ex]
  Table    & Lat  & Long  & Code  & Course/Speed &  \\ 
  \hline 
   &  &  &  & Compressed  & Type \\  [-1ex]
  
   &  &  &  & Radio Range &  \\ 
  \hline 
  Id & YYYY & XXXX &  & Compressed  & T \\ [-1ex]

      &  &  &  & Altitude &  \\
  \hline 
  1 & 4 & 4 & 1 & 2 & 1 \\
  \hline 
\end{tabular}



Compressed format can be used in place of lat/long position format
anywhere that

…ddmm.hhN/dddmm.hhW\$xxxxxxx…

occurs.

All bytes except for the / and \$ are base-91 printable ASCII
characters (!..\{). These are converted to numeric values by
subtracting 33 from the decimal ASCII character code. For example, \#
has an ASCII code of 35, and represents a numeric value of 2
(i.e. 35-33).

\section {Symbol}

The presence of the leading Symbol Table Identifier instead of a digit
indicates that this is a compressed Position Report and not a normal lat/long
report.


\section{Lat/Long Encoding}

The values of YYYY and XXXX are computed as follows:

YYYY is 380926 x (90 – latitude) [base 91]
latitude is positive for north, negative for south, in degrees.

XXXX is 190463 x (180 + longitude) [base 91]
longitude is positive for east, negative for west, in degrees.

For example, for a longitude of 72° 45' 00" west (i.e. -72.75
degrees), the math is 190463 x (180 – 72.75) = 20427156. Because this
is to base 91, it is then necessary to progressively divide this value
by reducing powers of 91, to obtain the numerical values of X:

20427156 / $91^{3}$ = 27, remainder 80739 \\
80739 / $91^{2} = 9$, remainder 6210 \\
6210 / $91^{1} = 68$, remainder 22 \\



To obtain the corresponding ASCII characters, 33 is added to each of
these values, yielding 60 (i.e. 27+33), 42, 101 and 55. From the ASCII
Code Table (in Appendix 3), this corresponds to <*e7 for XXXX.

\section{Lat/Long Decoding}

To decode a compressed lat/long, the reverse process is needed. That
is, if YYYY is represented as $y_{1}y_{2}y_{3}y_{4}$ and XXXX as
$x_{1}x_{2}x_{3}x_{4}$, then:


Lat $ = 90 - ((y_{1}-33) \times 91 + (y_{2}-33) \times 91 + (y_{3}-33) \times 91 + y_{4}-33) / 380926 $

Long $ = -180 + ((x1-33) \times 91 + (x2-33) x 91 + (x3-33) \times 91 + x4-33) / 190463 $

For example, if the compressed value of the longitude is <*e7 (as computed
above), the calculation becomes:


Long $ = -180 + (27 \times 91 + 9 \times 91 + 68 \times 91 + 22) / 190463 $

$ = -180 + (20346417 + 74529 + 6188 + 22) / 190463 $ \\
$ = -180 + 107.25 $ \\
$ = -72.75 $ degrees 

\section {Course/Speed, Pre-Calculated Radio Range and Altitude}

The two cs bytes following the Symbol Code character can contain either
the compressed course and speed or the compressed pre-calculated radio
range or the station’s altitude. These two bytes are in base 91 format.

In the special case of c = \textvisiblespace (space), there is no course, speed or range
data, in which case the csT bytes are ignored.

{\em Course/Speed} — If the ASCII code for c is in the range ! to z
inclusive — corresponding to numeric values in the range 0–89 decimal
(i.e. after subtracting 33 from the ASCII code) — then cs represents a
compressed course/speed value:


course $ = c \times 4 $

speed $ = 1.08^{s} – 1 $

For example, if the cs characters are 7P, the corresponding values of
c and s (after subtracting 33 from the ASCII character code) are 22
and 47 respectively. Substituting these values in the above equations:

course $ = 22 \times 4 = 88 $ degrees

speed $ = 1.0847 – 1 = 36.2 $ knots

{\em Pre-Calculated Radio Range} — If c = \{, then cs represents a
compressed pre-calculated radio range value:
  
range $ = 2 \times 1.08^{s} $

For example, if the cs bytes are {?, the ASCII code for ? is 63, so the value
of s is 30 (i.e. 63-33). Thus:

range $ = 2 \times 1.08^{30} \approx 20 $ miles

So APRS will draw a circle of radius 20 miles around the station plot on the
screen.

\section{The Compression Type (T) Byte}

The T byte follows the cs bytes. The T byte contains several bit fields
showing the GPS fix status, the NMEA source of the position data and the
origin of the compression.

The T byte is not meaningful if the c byte is \textvisiblespace (space).

% Compression Type (T) Byte Format Table

\begin{tabular}{c|c|c|c|c|c|}
  \hline
   \multicolumn {6}{r}{Compression Type (T) Byte Format} \\
  \hline
\end{tabular}

NMEA Source
0 0 = other
0 1 = GLL
1 0 = GGA
1 1 = RMC


Compression Origin
0 0 0 = Compressed
0 0 1 = TNC BText
0 1 0 = Software (DOS/Mac/Win/+SA)
0 1 1 = [tbd]
1 0 0 = KPC3
1 0 1 = Pico
1 1 0 = Other tracker [tbd]
1 1 1 = Digipeater conversion

For example, if the compressed position was derived from an RMC sentence,
the fix is current, and the compression was performed by APRSdos software,
then the value of T in binary is 0 0 1 11 010, which equates to 58 decimal.
Adding 33 to this value gives the ASCII code for the T byte (i.e. 91), which
corresponds to the [ character.

Thus, using data from all the earlier examples, if the RMC sentence contains
(among other parameters) the following data:

Latitude = 49° 30' 00" north
Longitude = 72° 45' 00" west
Speed = 36.2 knots
Course = 88°

and:


the fix is current
compression is performed by APRSdos software
the display symbol is a “car”

then the complete 13-character compressed location field is transmitted as:

% small table

\section{Altitude}


If the T byte indicates that the raw data originates from a GGA sentence (i.e.
bits 4 and 3 of the T byte are 10), then the sentence contains an altitude
value, in feet. After compression, the compressed altitude data is placed in
the cs bytes, such that:

altitude = 1.002cs feet

For example, if the received cs bytes are S], the computation is as follows:

% fill in




Subtract 33 from the ASCII code for each character:
c = 83 – 33 = 50
s = 93 – 33 = 60



Multiply c by 91 and add s to obtain cs:
cs = 50 x 91 + 60
= 4610



$ Then altitude = 1.002^{4610} $
= 10004 feet

\section {New Trackers}

Tracker firmware may compress GPS data directly to APRS compressed
format. They would use the ! Data Type Indicator, showing that the position
is real-time and that the tracker is not APRS-capable.
If the Position Report is not real-time, then the / Data Type Indicator can be
used instead, so that the latest fix time may be included.


Some digipeaters have the ability to convert raw NMEA strings from existing
trackers to compressed data format for further forwarding.

\section {Old Trackers}

These digipeaters will compress the data if the tracker Destination Address is
GPS. (Note: This is the 3-letter address GPS, not GPS*).
Trackers desiring for their packets to not be modified by the APRS network
will use any other valid generic APRS Destination Address.

Compressed data is contained in the AX.25 Information field, in these
formats:

\section {Compressed Report Formats}

% Table here

Examples
=/5L!!<*e7>V sTComment

2

with APRS messaging. Note the V space character following the >
Symbol Code, indicating that there is no course/speed, radio range or
altitude. The sT characters are fillers and have no significance here.
with APRS messaging, RMC sentence, with course/speed.
with APRS messaging, with radio range.
with APRS messaging, GGA sentence, altitude.



=/5L!!<*e7>7P[
=/5L!!<*e7>{?!
=/5L!!<*e7OS]S

% table here

Compressed Lat/Long Position Report Format — with Timestamp


Example
@092345z/5L!!<*e7>{?!


with APRS messaging, timestamp, radio range.


